% ===================================================================
%  ਸਾਰਣੀ ਨੰ. 9 — ਪਹਾੜ ਤੇ ਪਾਣੀਆਂ ਵਾਲਾ ਸ਼ਹਿਰ। — ਜੰਗਲ ਤੇ ਸ਼ਿਕਾਰ।
%  Traduction 2026 d'apres texto/io, controlee sur texto/fr, texto/en
%  et texto/hi. Consigne : voir texto/pa/10-sarni-01.tex.
% ===================================================================

%%K t09-apar-1 apar
\VUsaut{4.0mm}
\VUtitre{15.0pt}{60}{ਸਾਰਣੀ ਨੰ. 9}
\VUsaut{3.0mm}

%%K t09-tit-1 sub
\VUcentre{12.0pt}{0}{\textit{ਪਹਿਲਾ ਦ੍ਰਿਸ਼।}}
\VUsaut{3.0mm}

%%K t09-tit-2 sub
\VUpk{11.4pt}{40}{ਪਹਾੜ। --- ਪਾਣੀਆਂ ਵਾਲਾ ਸ਼ਹਿਰ।}
\VUsaut{3.0mm}

%%K t09-c1-01-1 p
1. --- ਮੈਂ ਇਕ ਹਫ਼ਤੇ ਤੋਂ ਪ੍ਰਾਤਜ਼ ਵਿਚ ਹਾਂ। ਪਿਰੇਨੀਜ਼ ਦੀ ਹੈਰਾਨ ਕਰ ਦੇਣ ਵਾਲੀ
\VUgras{ਪਹਾੜੀ ਲੜੀ} \textsuperscript{(1)} ਦੇ ਸਾਹਮਣੇ ਖੜੇ ਹੋ ਕੇ ਮੈਨੂੰ ਜੋ ਹੈਰਾਨੀ ਹੋਈ,
ਉਹ ਮੈਂ ਦੱਸ ਨਹੀਂ ਸਕਦਾ; ਉਸ ਦੀ \VUgras{ਚੋਟੀਆਂ ਦੀ ਲਕੀਰ} \textsuperscript{(2)} ਨੀਲੇ
ਅਸਮਾਨ ਉੱਤੇ ਕਿਸੇ ਵੱਡੀ ਝਾਲਰ ਵਾਂਗ ਉੱਭਰਦੀ ਹੈ।

%%K t09-c1-02-1 p
2. --- ਮੈਂ ਸੋਚਿਆ ਵੀ ਨਾ ਸੀ ਕਿ ਪਿਰੇਨੀਜ਼ ਦੀਆਂ \VUgras{ਸਿਖਰਾਂ} \textsuperscript{(3)}
ਇੰਨੀ \VUgras{ਉਚਾਈ} ਤਕ ਪਹੁੰਚ ਸਕਦੀਆਂ ਹਨ, ਤੇ ਮੈਂ ਉਨ੍ਹਾਂ ਸ਼ਾਨਦਾਰ
\VUgras{ਬਰਫ਼ਾਨੀ ਦਰਿਆਵਾਂ} \textsuperscript{(5)} ਤੇ ਬਰਫ਼ ਨਾਲ ਢਕੀਆਂ
\VUgras{ਚੋਟੀਆਂ} \textsuperscript{(4)} ਨੂੰ ਵੇਖ ਕੇ ਦੰਗ ਰਹਿ ਗਿਆ, ਜਿਨ੍ਹਾਂ ਤੋਂ ਇੰਨੇ
ਭਿਆਨਕ \VUgras{ਬਰਫ਼ ਦੇ ਤੋਦੇ} ਉਤਰਦੇ ਹਨ।

%%K t09-tit-3 sub
\VUsaut{3.0mm}
\VUpk{10.2pt}{40}{ਪਹਾੜੀ ਦ੍ਰਿਸ਼।}
\VUsaut{3.0mm}

%%K t09-c1-03-1 p
3. --- ਆਉਣ ਦੇ ਅਗਲੇ ਹੀ ਦਿਨ ਮੈਂ ਪ੍ਰਾਤਜ਼ ਕੋਲ ਉਸ ਪੁਰਾਣੇ ਬੁਝੇ \VUgras{ਅਗਨ-ਪਰਬਤ}
\textsuperscript{(6)} ਤਕ ਗਿਆ। \VUgras{ਮੂੰਹ} ਦੇ ਨੰਗੇ ਬੰਜਰ ਕੰਢੇ ਉੱਤੇ ਹੁਣ ਵੀ ਉਹ ਨਿਸ਼ਾਨ
ਸਾਫ਼ ਵਿਖਾਈ ਦਿੰਦੇ ਹਨ ਜੋ ਉਸ \VUgras{ਲਾਵੇ} ਨੇ ਛੱਡੇ ਸਨ, ਜੋ ਕਈ ਸਦੀਆਂ ਪਹਿਲਾਂ
\VUgras{ਪਹਾੜ ਦੀ ਢਾਲ} \textsuperscript{(7)} ਤੋਂ ਵਗਿਆ ਸੀ। ਇਕ \VUgras{ਢਲਾਣ} ਉੱਤੇ
ਮੈਨੂੰ ਉਸ ਲਾਵੇ ਦੇ ਕੁਝ ਟੁਕੜੇ ਮਿਲ ਹੀ ਗਏ।

%%K t09-c1-04-1 p
4. --- ਪ੍ਰਾਤਜ਼ ਹੱਦ ਉੱਤੇ ਵੱਸਿਆ ਹੈ, ਤੇ ਜੋ \VUgras{ਕਿਲ੍ਹਾ} \textsuperscript{(8)}
ਫ਼ਰਾਂਸ ਤੇ ਸਪੇਨ ਦੇ ਵਿਚਕਾਰਲੇ \VUgras{ਦੱਰੇ} \textsuperscript{(27)} ਦੀ ਰਾਖੀ ਕਰਦਾ ਹੈ,
ਉਹ ਰਾਈਨ ਦੇ ਕੰਢੇ ਦੇ ਉਨ੍ਹਾਂ ਪੁਰਾਣੇ ਕਿਲ੍ਹਿਆਂ ਵਰਗਾ ਹੀ ਹੈ ਜੋ ਆਪਣੀ ਚਟਾਨ ਦੇ ਉੱਤੋਂ ਕਿਸੇ
\VUgras{ਸ਼ਿਕਾਰ} ਦੀ ਤਾਕ ਵਿਚ ਲਗਦੇ ਹਨ।

%%K t09-c1-05-1 p
5. --- ਇਕ ਵੱਡਾ \VUgras{ਉਕਾਬ} \textsuperscript{(9)}, ਜਿਸ ਦਾ ਆਲ੍ਹਣਾ \VUgras{ਗੜ੍ਹੀ}
\textsuperscript{(8)} ਤੋਂ ਦੂਰ ਨਹੀਂ, ਅਕਸਰ ਮੇਰਾ ਧਿਆਨ ਖਿੱਚਦਾ ਹੈ। ਇਹ
\VUgras{ਸ਼ਿਕਾਰੀ ਪੰਛੀ}, ਆਪਣੀ ਤਕੜੀ \VUgras{ਚੁੰਝ} ਨਾਲ, ਅਕਸਰ ਆਪਣੇ ਬੱਚਿਆਂ ਲਈ ਕੋਈ
ਲੇਲਾ ਜਾਂ ਬੱਕਰੀ ਦਾ ਬੱਚਾ ਆਪਣੇ \VUgras{ਪੰਜਿਆਂ} ਵਿਚ ਦੱਬੀ ਲਿਆਉਂਦਾ ਹੈ। ਉਸ ਦੇ
\VUgras{ਖੰਭ} ਇੰਨੇ ਵੱਡੇ ਹਨ ਕਿ ਉਹ ਆਪਣੇ ਭਾਰੇ ਬੋਝ ਨਾਲ ਦੇਰ ਤਕ ਤਰ ਸਕਦਾ ਹੈ।

%%K t09-tit-4 sub
\VUsaut{3.0mm}
\VUpk{10.2pt}{40}{ਪੈਦਲ ਤੁਰਨ ਵਾਲਾ।}
\VUsaut{3.0mm}

%%K t09-c1-06-1 p
6. --- ਇੱਥੇ ਦੀ ਸਭ ਤੋਂ ਸੋਹਣੀ ਸੈਰ ਕੋ ਦੇ \VUgras{ਝਰਨੇ} \textsuperscript{(10)} ਤਕ
ਜਾਣਾ ਹੈ। \VUgras{ਡਿੱਗਦਾ ਪਾਣੀ} ਵਰੋਲਾ ਬਣ ਕੇ ਹੇਠਾਂ ਆਉਂਦਾ ਹੈ ਤੇ ਫਿਰ ਸ਼ਹਿਰ ਦੇ ਲਹਿੰਦੇ
ਦੇ \VUgras{ਦਿਆਰ ਦੇ ਜੰਗਲ} \textsuperscript{(11)} ਵਿਚ ਇਕ ਸੋਹਣੇ \VUgras{ਨਾਲੇ} ਦੇ ਰੂਪ
ਵਿਚ ਗੁੰਮ ਹੋ ਜਾਂਦਾ ਹੈ। ਅਸੀਂ ਇਸ ਜੰਗਲ ਵਿਚੋਂ \VUgras{ਮੌਸਮ ਦੇ ਕੇਂਦਰ}
\textsuperscript{(12)} ਤਕ ਚੜ੍ਹੇ, ਤੇ \VUgras{ਮੁਨਾਰੇ} ਦੇ ਉੱਤੋਂ ਅਸੀਂ ਇਲਾਕੇ ਦਾ ਸਾਰਾ
\VUgras{ਪੰਛੀ ਵਰਗਾ ਦ੍ਰਿਸ਼} ਸਮੇਟ ਲਿਆ। \VUgras{ਕੁਦਰਤੀ ਦ੍ਰਿਸ਼} ਬਹੁਤ ਵਧੀਆ ਹੈ, ਤੇ
ਲਗਦਾ ਹੈ \VUgras{ਕੁਦਰਤ} ਨੇ ਇਸ ਦੇਸ ਉੱਤੇ ਆਪਣੇ ਸਾਰੇ ਖ਼ਜ਼ਾਨੇ ਲੁਟਾ ਦਿੱਤੇ ਹਨ।

%%K t09-c1-07-1 p
7. --- ਹੇਠਾਂ ਉਤਰਨ ਲਈ ਅਸੀਂ \VUgras{ਰੱਸੇ ਵਾਲੀ ਰੇਲ} \textsuperscript{(13)} ਲਈ ਤੇ ਇਕ
ਨਿੱਕੇ \VUgras{ਸਟੇਸ਼ਨ} ਉੱਤੇ ਰੁਕੇ, ਇਕ ਪੁਰਾਣੇ \VUgras{ਜਾਗੀਰਦਾਰੀ ਕਿਲ੍ਹੇ} ਦੇ
\VUgras{ਖੰਡਰਾਂ} \textsuperscript{(14)} ਕੋਲ, ਜਿਸ ਵਿਚ ਕਦੇ ਪ੍ਰਾਤਜ਼ ਦੇ ਜਾਗੀਰਦਾਰ ਰਹਿੰਦੇ
ਸਨ।

%%K t09-c1-08-1 p
8. --- ਵਿਚਾਰਾ ਕਿਲ੍ਹਾ! ਉਹ ਕੀ ਸੋਚਦਾ ਹੋਵੇਗਾ ਜਦੋਂ ਆਪਣੇ ਹੇਠਾਂ ਉਸ ਸੁਥਰੇ
\VUgras{ਪਾਣੀਆਂ ਵਾਲੇ ਸ਼ਹਿਰ} \textsuperscript{(15)} ਨੂੰ ਵੇਖਦਾ ਹੈ, ਜੋ ਅੱਜ ਇਕ ਫ਼ੈਸ਼ਨ ਦਾ
\VUgras{ਗਰਮ ਚਸ਼ਮਿਆਂ ਦਾ ਸਿਹਤ-ਘਰ} ਹੈ?

%%K t09-c1-08-2 p
\VUgras{ਮੌਸਮ} ਇੰਨਾ ਸੋਹਣਾ ਹੈ ਤੇ \VUgras{ਖਣਿਜ ਪਾਣੀ} ਇੰਨਾ ਗੁਣਕਾਰੀ ਕਿ
\VUgras{ਸਿਹਤ-ਘਰ} \textsuperscript{(17)} ਵਿਚ ਲਿਆ \VUgras{ਇਲਾਜ ਦਾ ਦੌਰ} ਸੱਚਮੁੱਚ ਸੁਖ
ਦੀ ਗੱਲ ਹੈ।

%%K t09-tit-5 sub
\VUsaut{3.0mm}
\VUpk{10.2pt}{40}{ਇਕ ਸੋਹਣਾ ਪਾਣੀਆਂ ਵਾਲਾ ਸ਼ਹਿਰ।}
\VUsaut{3.0mm}

%%K t09-c1-09-1 p
9. --- ਸੱਚਮੁੱਚ, ਟਿਕਾਅ ਨੂੰ ਸੋਹਣਾ ਬਣਾਉਣ ਲਈ ਸਭ ਕੁਝ ਕੀਤਾ ਗਿਆ ਹੈ। ਰੇਲ ਲੈ ਜਾਣ ਲਈ ਇਕ
ਵੱਡਾ \VUgras{ਪੁਲ} \textsuperscript{(16)} ਬਣਾਇਆ ਗਿਆ; \VUgras{ਪਾਣੀਆਂ ਦੇ ਘਰ} ਦਾ
\VUgras{ਬਾਗ਼} \textsuperscript{(18)}, ਜਿੱਥੇ \VUgras{ਸੰਗੀਤ ਦੀਆਂ ਮਹਿਫ਼ਲਾਂ} ਹੁੰਦੀਆਂ
ਹਨ, ਬੜੇ ਸਲੀਕੇ ਨਾਲ ਸਜਾਇਆ ਗਿਆ ਹੈ। ਕਈ ਸੋਹਣੀਆਂ \VUgras{ਕੋਠੀਆਂ}
\textsuperscript{(19)} \VUgras{ਜੂਆ-ਘਰ} \textsuperscript{(21)} ਦੇ ਚੁਫੇਰੇ ਬਣੀਆਂ ਹਨ, ਤੇ
ਇਕ ਬਹੁਤ ਚੰਗੀ ਤਰ੍ਹਾਂ ਰੱਖੀ \VUgras{ਸ਼ਾਹ-ਰਾਹ} \textsuperscript{(22)} ਆਉਣਾ-ਜਾਣਾ ਬਹੁਤ
ਸੌਖਾ ਕਰ ਦਿੰਦੀ ਹੈ।

%%K t09-c1-10-1 p
10. --- ਇਕ ਸਾਦਾ \VUgras{ਬੰਨ੍ਹ} \textsuperscript{(23)} \VUgras{ਸੜਕ}
\textsuperscript{(20)} ਨੂੰ ਅਸਲੀ \VUgras{ਵਾਦੀ} \textsuperscript{(25)} ਤੋਂ ਵੱਖ ਕਰਦਾ ਹੈ,
ਜਿਸ ਦੇ ਤਲੇ ਵਿਚ ਇਕ ਸੋਹਣੀ ਨਿੱਕੀ \VUgras{ਝੀਲ} \textsuperscript{(26)} ਹੈ, ਜੋ ਇਕ ਸਾਫ਼
\VUgras{ਨਾਲੇ} \textsuperscript{(24)} ਨੂੰ ਪਾਣੀ ਦਿੰਦੀ ਹੈ।

%%K t09-c1-11-1 p
11. --- ਵਾਦੀ ਦੇ ਦੂਜੇ ਪਾਸੇ ਲੋਹੇ ਦੀ ਇਕ \VUgras{ਕਾਨ} \textsuperscript{(28)} ਪੁੱਟੀ ਜਾ
ਰਹੀ ਹੈ। ਕਈ ਵਾਰ ਮੈਂ \VUgras{ਕਾਨ-ਕਾਮਿਆਂ} ਨੂੰ \VUgras{ਖੂਹ} ਵਿਚ ਉਤਰਦੇ ਵੇਖਣ ਗਿਆ, ਤੇ
ਮੈਂ ਉਸ \VUgras{ਸੁਰੰਗ} ਵਿਚ ਵੀ ਗਿਆ ਜਿੱਥੇ ਉਹ ਸਾਰਾ ਦਿਨ ਉਹ \VUgras{ਕੱਚੀ ਧਾਤ}
ਕੱਢਦੇ ਹਨ ਜਿਸ ਵਿਚ \VUgras{ਧਾਤ} ਰਹਿੰਦੀ ਹੈ।

%%K t09-c1-12-1 p
12. --- ਕਾਨ ਕੋਲ ਇਕ ਵੱਡੀ \VUgras{ਢਲਾਈ-ਸ਼ਾਲਾ} ਬਣਾਈ ਗਈ, ਤਾਂ ਜੋ ਕੱਢੀ ਹੋਈ ਦੌਲਤ ਉੱਥੇ
ਹੀ ਕੰਮ ਆ ਜਾਵੇ। \VUgras{ਭੱਠੀ} \textsuperscript{(29)}, ਜੋ ਇੱਥੇ ਬਹੁਤਾਤ ਨਾਲ ਲੱਭਦੇ
\VUgras{ਕੋਲੇ} ਨਾਲ ਗਰਮ ਕੀਤੀ ਜਾਂਦੀ ਹੈ, ਛੇਤੀ ਤੇ ਸਸਤੇ ਵਿਚ \VUgras{ਕੱਚਾ ਲੋਹਾ} ਲੱਭਣ
ਦਿੰਦੀ ਹੈ।

%%K t09-c1-13-1 p
13. --- ਕਾਨ ਤੋਂ ਦੂਰ ਨਹੀਂ ਇਕ \VUgras{ਖਾਨ} \textsuperscript{(30)} ਕੱਟੀ ਜਾ ਰਹੀ ਹੈ।
ਬੁੱਢਾ \VUgras{ਖਾਨ ਦਾ ਮਜ਼ਦੂਰ} ਆਪਣੀ \VUgras{ਕਹੀ} ਨਾਲ ਜੋ ਤੋੜਦਾ ਹੈ ਉਹ ਨਾ
\VUgras{ਗਰੇਨਾਈਟ} ਹੈ, ਨਾ \VUgras{ਪੋਰਫ਼ਰੀ}, ਨਾ \VUgras{ਰੇਤਲਾ ਪੱਥਰ}, ਸਗੋਂ ਘਰ ਬਣਾਉਣ
ਦਾ ਸਧਾਰਨ \VUgras{ਇਮਾਰਤੀ ਪੱਥਰ}।

%%K t09-c1-14-1 p
14. --- ਇਸ ਦੇਸ ਦੇ ਸਭ ਤੋਂ ਸੋਹਣੇ ਗਹਿਣਿਆਂ ਵਿਚੋਂ ਇਕ ਹੈ \VUgras{ਦਿਆਰ}
\textsuperscript{(31)}। ਹਰ ਥਾਂ — ਕਿਸੇ \VUgras{ਖੱਡ} \textsuperscript{(32)} ਦੇ ਤਲੇ
ਵਿਚ, ਕਿਸੇ \VUgras{ਆਬਸ਼ਾਰ} \textsuperscript{(33)} ਦੇ ਕੰਢੇ, ਕਿਸੇ \VUgras{ਗੁਫ਼ਾ}
\textsuperscript{(34)} ਜਾਂ ਕਗਾਰ ਕੋਲ — ਉਸ ਦੀਆਂ ਪਤਲੀਆਂ ਸੂਈਆਂ ਆਪਣਾ ਹਰਾ ਸੁਰ ਜੋੜ
ਦਿੰਦੀਆਂ ਹਨ।

%%K t09-tit-6 sub
\VUsaut{3.0mm}
\VUpk{10.2pt}{40}{ਪੈਦਲ ਤੁਰਨਾ।}
\VUsaut{3.0mm}

%%K t09-c1-15-1 p
15. --- ਮੈਂ ਆਪਣੀਆਂ ਛੁੱਟੀਆਂ ਦਾ ਫ਼ਾਇਦਾ ਲੰਮੀਆਂ ਸੈਰਾਂ ਲਈ ਚੁੱਕ ਰਿਹਾ ਹਾਂ। ਜੋ
\VUgras{ਪਗਡੰਡੀਆਂ} \textsuperscript{(35)} ਪਹਾੜ ਉੱਤੇ ਚੱਕਰ ਕੱਟਦੀਆਂ ਚੜ੍ਹਦੀਆਂ ਹਨ,
ਉਨ੍ਹਾਂ ਉੱਤੇ ਕਈ \VUgras{ਕਿਲੋਮੀਟਰ}, ਤੇ ਅਕਸਰ ਕਈ \VUgras{ਕੋਹ}, ਪੈਂਡਾ ਪਤਾ ਲੱਗੇ ਬਿਨਾਂ
ਮੁੱਕ ਜਾਂਦਾ ਹੈ।

%%K t09-c1-15-2 p
\VUgras{ਮੀਲ-ਪੱਥਰ} \textsuperscript{(36)} ਤੇ \VUgras{ਸੇਧ-ਸੂਚਕ}
\textsuperscript{(37)} ਰਾਹਾਂ ਨੂੰ ਨਿਸ਼ਾਨ ਲਾਉਂਦੇ ਹਨ, ਜਿਨ੍ਹਾਂ ਉੱਤੇ ਅਕਸਰ ਕੋਈ ਜਵਾਨ
\VUgras{ਪਹਾੜੀ ਮੁੰਡਾ} \textsuperscript{(38)} ਮਿਲਦਾ ਹੈ, ਕੱਛ ਵਿਚ \VUgras{ਝੋਲਾ}
\textsuperscript{(40)} ਚੁੱਕੀ ਤੇ ਇਕ \VUgras{ਮਾਰਮੋਟ} \textsuperscript{(39)} ਚੁੱਕਿਆ,
ਜਾਂ ਕੋਈ \VUgras{ਬੰਜਾਰਾ} \textsuperscript{(41)}, ਜਿਸ ਦੀ \VUgras{ਜੱਤ ਵਾਲੀ ਟੋਪੀ}
\textsuperscript{(42)} ਉਸ ਦੇ ਪਾਟੇ ਪੁਰਾਣੇ \VUgras{ਓੜ੍ਹਨ} \textsuperscript{(43)} ਨਾਲ
ਅਜੀਬ ਮੇਲ ਖਾਂਦੀ ਹੈ। ਉਹ ਇਕ ਸਿਧਾਏ \VUgras{ਰਿੱਛ} \textsuperscript{(44)} ਨੂੰ
\VUgras{ਸੰਗਲ} ਨਾਲ ਲਈ ਤੁਰਦਾ ਹੈ।

%%K t09-tit-7 sub
\VUpk{10.2pt}{40}{ਚੜ੍ਹਾਈ।}
\VUsaut{3.0mm}

%%K t09-c1-16-1 p
16. --- ਤਕਰੀਬਨ ਹਰ ਦਿਨ ਮੈਂ ਇਕ \VUgras{ਗਾਈਡ} \textsuperscript{(48)} ਨਾਲ
\VUgras{ਚੜ੍ਹਾਈ} ਕਰਨ ਜਾਂਦਾ ਹਾਂ, ਤੇ ਮੈਨੂੰ ਬੜਾ ਮਾਣ ਹੁੰਦਾ ਹੈ ਜਦੋਂ ਪਿੱਠ ਉੱਤੇ ਆਪਣਾ
\VUgras{ਪਿੱਠੂ} \textsuperscript{(47)} ਤੇ ਹੱਥ ਵਿਚ ਆਪਣੀ \VUgras{ਲਾਠੀ} ਚੁੱਕੀ, ਕਿਸੇ
ਸੱਚੇ \VUgras{ਸੈਲਾਨੀ} \textsuperscript{(46)} ਵਾਂਗ, ਮੈਂ \VUgras{ਚਟਾਨਾਂ}
\textsuperscript{(45)} ਉੱਤੇ ਚੜ੍ਹਾਈ ਕਰਦਾ ਹਾਂ।

%%K t09-c1-17-1 p
17. --- ਪਹਾੜ ਦੀ ਚੋਟੀ ਤੋਂ ਮੈਦਾਨ ਦੇ ਲੋਕ ਕੀੜੀਆਂ ਤੋਂ ਵੱਡੇ ਨਹੀਂ ਲਗਦੇ;
\VUgras{ਭੇਡਾਂ ਦਾ ਇੱਜੜ} \textsuperscript{(50)}, ਜੋ \VUgras{ਚਰਾਗਾਹ}
\textsuperscript{(49)} ਵਿਚ ਚਰਦਾ ਹੈ, ਬੱਚਿਆਂ ਦਾ ਖਿਡੌਣਾ ਲਗਦਾ ਹੈ। ਇਕ \VUgras{ਚੀੜ੍ਹ}
\textsuperscript{(51)}, ਜਾਂ \VUgras{ਲੱਕੜ ਦਾ ਪਹਾੜੀ ਘਰ} \textsuperscript{(52)} ਤਕ, ਹਰੇ
ਉੱਤੇ ਇਕ ਬਿੰਦੂ ਤੋਂ ਵੱਧ ਨਹੀਂ।

%%K t09-c1-18-1 p
18. --- ਇਨ੍ਹਾਂ ਸੈਰਾਂ ਵਿਚ ਸਾਨੂੰ \VUgras{ਪਗਡੰਡੀ} \textsuperscript{(53)} ਉੱਤੇ ਅਕਸਰ
ਕੋਈ \VUgras{ਸ਼ਾਮੋਆ ਦਾ ਸ਼ਿਕਾਰੀ} \textsuperscript{(54)} ਮਿਲਦਾ ਹੈ, ਜੋ ਹੁਣੇ ਹੁਣੇ ਮਾਰਿਆ
ਹੋਇਆ ਡੰਗਰ ਮੋਢੇ ਉੱਤੇ ਚੁੱਕੀ ਹੁੰਦਾ ਹੈ।

%%K t09-c1-19-1 p
19. --- ਮੈਂ ਆਪਣੇ ਬੂਟਿਆਂ ਦੇ ਸੰਗ੍ਰਹਿ ਵਿਚ \VUgras{ਫ਼ਰਨ} \textsuperscript{(55)} ਦਾ ਇਕ
ਬਹੁਤ ਅਨੋਖਾ ਪੱਤਾ ਤੇ \VUgras{ਹੀਦਰ} \textsuperscript{(57)} ਦੀ ਇਕ ਸੋਹਣੀ ਗੁਲਾਬੀ ਟਾਹਣੀ
ਰੱਖੀ ਹੈ, ਜੋ ਮੈਂ ਪਿਛਲੀ ਸੈਰ ਵਿਚ ਇਕ ਨਿੱਕੀ \VUgras{ਗੁਫ਼ਾ} \textsuperscript{(56)} ਦੇ ਮੂੰਹ
ਉੱਤੇ ਇਕੱਠੀ ਕੀਤੀ ਸੀ।

%%K t09-tit-8 sub
\VUsaut{3.0mm}
\VUcentre{12.0pt}{0}{\textit{ਦੂਜਾ ਦ੍ਰਿਸ਼।}}
\VUsaut{3.0mm}

%%K t09-tit-9 sub
\VUpk{11.4pt}{40}{ਜੰਗਲ। --- ਸ਼ਿਕਾਰ।}
\VUsaut{3.0mm}

%%K t09-c2-01-1 p
1. --- ਕੱਲ੍ਹ ਮੈਂ ਜੰਗਲ ਵਿਚ ਇਕ ਆਨੰਦ ਭਰਿਆ ਦਿਨ ਲੰਘਾਇਆ। ਉੱਥੇ ਰਹਿਣ ਵਾਲੇ ਡੰਗਰਾਂ ਤੇ
\VUgras{ਕੀੜਿਆਂ} \textsuperscript{(5)} ਦੇ ਰੁੱਝੇ ਜੀਵਨ ਨੇ ਮੈਨੂੰ ਬਹੁਤ ਚੰਗਾ ਲੱਗਾ।

%%K t09-c2-01-2 p
\VUgras{ਮੱਕੜੀ} \textsuperscript{(1)}, ਜੋ ਦੋ ਟਾਹਣੀਆਂ ਵਿਚਕਾਰ ਆਪਣਾ \VUgras{ਜਾਲਾ}
\textsuperscript{(2)} ਬੁਣ ਕੇ ਲੰਘਦੀ \VUgras{ਮੱਖੀ} \textsuperscript{(3)} ਫੜਦੀ ਹੈ,
\VUgras{ਘੋਗਾ} \textsuperscript{(4)}, ਜੋ ਆਪਣਾ ਖੋਲ ਚੁੱਕੀ ਔਖਾ ਹੀ ਅੱਗੇ ਵਧਦਾ ਹੈ, ਉਹ
ਭੌਂਰਾ ਜੋ ਆਪਣੇ ਸ਼ਿਕਾਰ ਦੀ ਟੋਹ ਵਿਚ ਹੈ, ਉਹ ਮਿਹਨਤੀ ਕੀੜੀ ਜੋ \VUgras{ਭੌਣ}
\textsuperscript{(6)} ਕੋਲ ਕੰਮ ਵਿਚ ਜੁਟੀ ਹੈ — ਇਹ ਸਾਰੇ ਨਿੱਕੇ ਜੀਵ ਅਸਾਧਾਰਨ ਜੋਸ਼ ਨਾਲ
ਆਉਂਦੇ-ਜਾਂਦੇ ਤੇ ਕੰਮ ਕਰਦੇ ਰਹੇ।

%%K t09-c2-02-1 p
2. --- ਇਕ \VUgras{ਸੱਪ} \textsuperscript{(7)}, ਜਿਸ ਨੂੰ ਪਹਿਲਾਂ ਮੈਂ \VUgras{ਨਾਗ}
ਸਮਝਿਆ ਪਰ ਜੋ ਸਿਰਫ਼ ਇਕ ਬੇਕਸੂਰ \VUgras{ਪਾਣੀ ਦਾ ਸੱਪ} ਸੀ, ਇਕ ਮੁੱਢ ਦੇ ਹੇਠਾਂ ਕੁੰਡਲੀ ਮਾਰੇ
ਸੀ, ਤੇ ਵਿਚ ਵਿਚ ਆਪਣਾ ਪਤਲਾ ਸਿਰ ਕੱਢਦਾ ਸੀ, ਜੋ ਹਮੇਸ਼ਾ ਡੰਗਣ ਨੂੰ ਤਿਆਰ ਲਗਦਾ ਸੀ। ਮੇਰੇ
ਸਾਹਮਣੇ ਇਕ ਵੱਡਾ \VUgras{ਸਲੱਗ} \textsuperscript{(8)} ਭਾਰੀ ਚਾਲ ਨਾਲ ਰਿੜ੍ਹ ਰਿਹਾ ਸੀ,
ਉਨ੍ਹਾਂ ਸੁਆਦੀ ਨਿੱਕੀਆਂ \VUgras{ਜੰਗਲੀ ਸਟ੍ਰਾਬੇਰੀਆਂ} \textsuperscript{(11)} ਵਿਚੋਂ ਇਕ
ਚੱਟ ਕਰ ਲੈਣ ਪਿੱਛੋਂ, ਜੋ ਇੱਥੇ ਬਹੁਤਾਤ ਨਾਲ ਉੱਗਦੀਆਂ ਹਨ।

%%K t09-c2-03-1 p
3. --- ਜ਼ਮੀਨ \VUgras{ਜੰਗਲੀ ਫੁੱਲਾਂ} ਨਾਲ ਵਿਛੀ ਸੀ : \VUgras{ਪ੍ਰਿਮਰੋਜ਼}
\textsuperscript{(9)}, \VUgras{ਪੈਰੀਵਿੰਕਲ} \textsuperscript{(10)} ਤੇ ਹੋਰ ਵੀ ਬਹੁਤ ਸਾਰੇ
ਬੂਟੇ, ਜਿਨ੍ਹਾਂ ਨੂੰ ਮੈਂ ਜਾਣਦਾ ਨਾ ਸੀ, \VUgras{ਘਾਹ ਦੇ ਗੁੱਛਿਆਂ} \textsuperscript{(12)}
ਵਿਚਕਾਰ ਖਿੜੇ ਸਨ।

%%K t09-c2-04-1 p
4. --- ਇਸ ਇਲਾਕੇ ਵਿਚ \VUgras{ਟਰੱਫ਼ਲ} ਤਕਰੀਬਨ ਓਨਾ ਹੀ ਆਮ ਹੈ ਜਿੰਨੀ \VUgras{ਖੁੰਬ}
\textsuperscript{(14)}, ਤੇ ਇਕ ਜੰਗਲ ਦੇ ਰਾਖੇ ਦਾ \VUgras{ਸੂਰ ਦਾ ਬੱਚਾ}
\textsuperscript{(13)} ਲਾਲਚ ਨਾਲ ਲੱਭਣ ਵਿਚ ਲੱਗਾ ਸੀ ਕਿ ਕੁਝ ਲੱਭ ਜਾਵੇ।

%%K t09-tit-10 sub
\VUsaut{3.0mm}
\VUpk{10.2pt}{40}{ਚੋਰੀ ਦਾ ਸ਼ਿਕਾਰ।}
\VUsaut{3.0mm}

%%K t09-c2-05-1 p
5. --- ਮੈਂ ਇਕ ਦ੍ਰਿਸ਼ ਦਾ \VUgras{ਅੰਤ} ਵੇਖਿਆ ਜਿਸ ਨੇ ਮੈਨੂੰ ਬਹੁਤ ਹਸਾਇਆ। ਆਪਣੇ
\VUgras{ਖੋਜੀ ਕੁੱਤੇ} \textsuperscript{(15)} ਦੇ ਨੱਕ ਦੀ ਬਦੌਲਤ \VUgras{ਜੰਗਲ ਦੇ ਰਾਖੇ}
\textsuperscript{(16)} ਨੇ ਹੁਣੇ ਹੁਣੇ ਇਕ \VUgras{ਚੋਰ-ਸ਼ਿਕਾਰੀ} \textsuperscript{(22)}
ਨੂੰ ਉਸ ਵੇਲੇ ਰੰਗੇ ਹੱਥੀਂ ਫੜਿਆ ਸੀ ਜਦੋਂ ਉਹ ਆਪਣਾ ਮਾਰਿਆ \VUgras{ਸ਼ਿਕਾਰ}
\textsuperscript{(17)} ਚੁੱਕ ਲੈ ਜਾਣ ਦੀ ਤਿਆਰੀ ਵਿਚ ਸੀ। ਫੜੇ ਜਾਣ ਨਾਲੋਂ ਚੰਗਾ ਉਸ ਨੂੰ ਆਪਣਾ
ਸ਼ਿਕਾਰ ਛੱਡ ਦੇਣਾ ਲੱਗਾ, ਸੋ ਚੋਰ-ਸ਼ਿਕਾਰੀ ਨੱਸ ਗਿਆ, ਤੇ \VUgras{ਸਹਾ}
\textsuperscript{(18)}, \VUgras{ਹਿਰਨੀ} \textsuperscript{(19)},
\VUgras{ਬਾਰਾਂਸਿੰਗਾ} \textsuperscript{(20)} ਤੇ \VUgras{ਜੰਗਲੀ ਸੂਰ}
\textsuperscript{(21)} ਘਾਹ ਉੱਤੇ ਪਏ ਰਹਿ ਗਏ, ਜੰਗਲ ਦੇ ਰਾਖੇ ਦੀ ਖ਼ੁਸ਼ੀ ਲਈ।

%%K t09-c2-06-1 p
6. --- ਜੰਗਲ ਵਿਚ ਰੁੱਖਾਂ ਦੀਆਂ ਜਿੰਨੀਆਂ ਕਿਸਮਾਂ ਹਨ, ਉਹ ਦੰਗ ਕਰ ਦਿੰਦੀਆਂ ਹਨ।
\VUgras{ਬੀਚ} \textsuperscript{(23)} ਤੇ \VUgras{ਭੋਜਪੱਤਰ} \textsuperscript{(26)},
\VUgras{ਚੀੜ੍ਹ}, ਜਿਨ੍ਹਾਂ ਤੋਂ \VUgras{ਗੂੰਦ} \textsuperscript{(27)} ਲੱਭਦੀ ਹੈ,
\VUgras{ਬਲੂਤ} \textsuperscript{(29)} ਤੇ ਹੋਰ ਵੀ ਬਹੁਤ ਸਾਰੇ ਆਪਣੀਆਂ ਟਾਹਣੀਆਂ ਮਿਲਾਉਂਦੇ
ਹਨ।

%%K t09-c2-06-2 p
ਉੱਚੇ ਰੁੱਖਾਂ ਦੇ ਮੁੱਢਾਂ ਨਾਲ ਲਿਪਟੀ \VUgras{ਆਈਵੀ} \textsuperscript{(24)}
\VUgras{ਉੱਚੇ ਜੰਗਲ} \textsuperscript{(25)} ਨੂੰ ਇਕ ਚਮਕਦਾ ਹਰਾ ਰੰਗ ਦਿੰਦੀ ਹੈ, ਜੋ ਉਸ
\VUgras{ਕਾਈ} \textsuperscript{(31)} ਦੀ ਫਿੱਕੀ ਕੋਮਲ ਛਾਂ ਨਾਲ ਸੋਹਣੇ ਢੰਗ ਨਾਲ ਮਿਲ ਜਾਂਦਾ
ਹੈ ਜਿਸ ਵਿਚ \VUgras{ਹਰਾ ਠਕ-ਠਕ ਪੰਛੀ} \textsuperscript{(30)} ਕੀੜੇ ਲੱਭਦਾ ਹੈ।

%%K t09-tit-11 sub
\VUsaut{3.0mm}
\VUpk{10.2pt}{40}{ਜੰਗਲ ਦਾ ਦ੍ਰਿਸ਼।}
\VUsaut{3.0mm}

%%K t09-c2-07-1 p
7. --- ਪੁਰਾਣੇ ਬਲੂਤ ਮੈਨੂੰ ਮੋਹ ਲੈਂਦੇ ਹਨ। ਉਨ੍ਹਾਂ ਦੀਆਂ ਵੱਡੀਆਂ ਵਲ ਖਾਂਦੀਆਂ
\VUgras{ਜੜ੍ਹਾਂ} \textsuperscript{(32)}, ਅੱਧ ਜ਼ਮੀਨ ਤੋਂ ਬਾਹਰ, ਉਨ੍ਹਾਂ ਨੂੰ ਜ਼ੋਰ ਤੇ ਤੇਜ
ਦਾ ਸ਼ਾਨਦਾਰ ਰੂਪ ਦਿੰਦੀਆਂ ਹਨ, ਤੇ ਮੈਨੂੰ ਅਚਰਜ ਹੁੰਦਾ ਹੈ ਕਿ \VUgras{ਮਾਲਕ}
\textsuperscript{(35)} ਉਨ੍ਹਾਂ ਨੂੰ ਵਢਵਾ ਕੇ ਉਸ \VUgras{ਆਰੀ} \textsuperscript{(34)} ਨੂੰ
ਸੌਂਪਣ ਉੱਤੇ ਕਿਵੇਂ ਰਾਜ਼ੀ ਹੋ ਜਾਂਦਾ ਹੈ ਜੋ \VUgras{ਲੱਕੜਹਾਰੇ} \textsuperscript{(33)} ਦੀ
ਹੈ। ਵਿਚਾਰੇ \VUgras{ਮੁੱਢ} \textsuperscript{(36)}, ਆਪਣੀ \VUgras{ਛਿੱਲ}
\textsuperscript{(37)} ਤੋਂ ਛਿੱਲੇ ਜਾਂ \VUgras{ਕੁਹਾੜੇ} \textsuperscript{(38)} ਨਾਲ ਚੀਰੇ,
ਜੋ \VUgras{ਦਿਹਾੜੀਦਾਰ ਮਜ਼ਦੂਰ} \textsuperscript{(39)} ਦਾ ਹੈ, ਮੈਨੂੰ ਦੁੱਖ ਦਿੰਦੇ ਹਨ।

%%K t09-c2-08-1 p
8. --- ਕੋਲ ਹੀ ਇਕ ਬੁੱਢੇ \VUgras{ਕੋਲਾ ਬਣਾਉਣ ਵਾਲੇ} \textsuperscript{(41)} ਨੇ ਆਪਣੀ
\VUgras{ਕਹੀ} ਨਾਲ ਕੋਲੇ ਦਾ ਇਕ \VUgras{ਢੇਰ} \textsuperscript{(42)} ਲਾਇਆ ਹੋਇਆ ਸੀ, ਤੇ ਇਕ
ਬੁੱਢੀ ਵੱਢੇ ਰੁੱਖਾਂ ਦੀਆਂ ਟਾਹਣੀਆਂ ਤੋਂ ਬਣੀ \VUgras{ਸੁੱਕੀ ਲੱਕੜ} ਦੀ ਇਕ
\VUgras{ਗੱਠੜੀ} \textsuperscript{(40)} ਲੈ ਜਾ ਰਹੀ ਸੀ।

%%K t09-tit-12 sub
\VUsaut{3.0mm}
\VUpk{10.2pt}{40}{ਸ਼ਿਕਾਰ।}
\VUsaut{3.0mm}

%%K t09-c2-09-1 p
9. --- ਮੈਂ \VUgras{ਜੰਗਲ ਦੇ ਰਾਖੇ ਦੇ ਘਰ} \textsuperscript{(46)} ਵਿਚ ਥੋੜਾ ਸਾਹ ਲੈਣ
ਜਾਣ ਵਾਲਾ ਸੀ, ਪਰ ਉਸ ਨਿੱਕੇ \VUgras{ਤਾਲ} \textsuperscript{(43)} ਤਕ ਪਹੁੰਚ ਕੇ, ਜਿਸ ਉੱਤੇ
ਇਕ ਸੋਹਣਾ \VUgras{ਹੰਸ} \textsuperscript{(45)} ਤਰ ਰਿਹਾ ਸੀ, ਮੈਂ ਵੇਖਿਆ ਕਿ ਹਾਲ ਦੇ ਇਕ
\VUgras{ਹੜ੍ਹ} \textsuperscript{(44)} ਨੇ ਰਾਹ ਨੂੰ ਸੱਚਮੁੱਚ ਦਾ \VUgras{ਚਿੱਕੜ} ਬਣਾ
ਦਿੱਤਾ ਹੈ। ਮੈਨੂੰ ਚੱਕਰ ਕੱਟਣਾ ਪਿਆ, ਤੇ \VUgras{ਦਿਆਰ} \textsuperscript{(49)},
\VUgras{ਸਫ਼ੈਦਿਆਂ} \textsuperscript{(48)} ਤੇ \VUgras{ਐਲਮ} \textsuperscript{(47)} ਕੋਲੋਂ
ਲੰਘਦੇ ਹੋਏ, ਜੋ ਜੰਗਲ ਦੇ ਕੰਢੇ ਖੜੇ ਹਨ, ਮੈਂ ਵੇਖਿਆ ਕਿ ਇਕ \VUgras{ਝਾੜੀ}
\textsuperscript{(54)} ਵਿਚੋਂ ਨਿਕਲ ਕੇ ਇਕ ਸ਼ਾਨਦਾਰ \VUgras{ਬਾਰਾਂਸਿੰਗਾ}
\textsuperscript{(50)} ਨੱਸਿਆ, ਜਿਸ ਦੇ ਪਿੱਛੇ ਇਕ ਨਾ ਥੱਕਣ ਵਾਲਾ
\VUgras{ਸ਼ਿਕਾਰੀ ਕੁੱਤਿਆਂ ਦਾ ਝੁੰਡ} \textsuperscript{(51)} ਲੱਗਾ ਸੀ, ਜਿਸ ਨੂੰ
\VUgras{ਨਰਸਿੰਘਾ} \textsuperscript{(53)} ਉਕਸਾ ਰਿਹਾ ਸੀ,
\VUgras{ਘੋੜਿਆਂ ਉੱਤੇ ਸਵਾਰ ਸ਼ਿਕਾਰੀਆਂ} \textsuperscript{(52)} ਦਾ। ਵਿਚਾਰਾ ਡੰਗਰ ਬਿਲਕੁਲ
ਥੱਕ ਚੁੱਕਾ ਸੀ। ਉਹ ਮੇਰੇ ਤੋਂ ਕੁਝ ਹੀ ਕਦਮਾਂ ਉੱਤੇ ਮਰਦਾ ਹੋਇਆ ਡਿੱਗ ਪਿਆ।

%%K t09-c2-10-1 p
10. --- ਜੰਗਲ ਦੇ ਰਾਖੇ ਦਾ ਪੁੱਤਰ, ਜਿਸ ਨੂੰ \VUgras{ਸ਼ਿਕਾਰ} ਦੇ ਰੌਲੇ ਨੇ ਬਾਹਰ ਖਿੱਚ
ਲਿਆ, ਘਰੋਂ ਨਿਕਲਿਆ, ਤੇ ਅਸੀਂ ਦੋਵੇਂ ਮੁੜ ਜੰਗਲ ਵਿਚ ਆਏ। ਉਸ ਨੇ ਇਕ ਪੁਰਾਣੇ ਬਲੂਤ ਦੀ ਟਾਹਣੀ
ਉੱਤੇ ਇਕ \VUgras{ਨੀਲਕੰਠ} \textsuperscript{(56)} ਵੇਖਿਆ ਸੀ। ਉਸ ਨੂੰ ਲੱਗਾ ਕਿ ਉਸ ਦਾ ਆਲ੍ਹਣਾ
\VUgras{ਪੱਤਿਆਂ} \textsuperscript{(58)} ਵਿਚ ਲੁਕਿਆ ਹੋਵੇਗਾ, ਤੇ ਉਹ ਉਸ ਨੂੰ ਲੈਣਾ ਚਾਹੁੰਦਾ
ਸੀ।

%%K t09-c2-10-2 p
\VUgras{ਗਾਲ੍ਹੜ} \textsuperscript{(61)} ਵਾਂਗ ਫੁਰਤੀਲਾ, ਉਹ ਇਕ \VUgras{ਟਾਹਣੀ}
\textsuperscript{(55)} ਤੋਂ ਦੂਜੀ ਟਾਹਣੀ ਚੜ੍ਹਦਾ ਗਿਆ, ਪਰ ਉਸ ਨੂੰ ਕੁਝ ਨਾ ਲੱਭਿਆ ਤੇ ਉਹ
ਸਿਰਫ਼ ਇਕ ਬੁੱਢੇ \VUgras{ਕਾਂ} \textsuperscript{(60)} ਨੂੰ ਚੌਂਕਾ ਸਕਿਆ, ਜੋ ਇਕ
\VUgras{ਸ਼ਾਖ਼} \textsuperscript{(57)} ਉੱਤੇ ਬੈਠਾ ਸੀ।
\VUsaut{4.0mm}
\VUfilet{22mm}
